\section{Conclusion and Future Work}
\label{sec:conclusion}

We set out to study the recall--memory tradeoff in subquadratic sequence models with maximal
honesty, and report three kinds of outcome. \emph{Theory:} we prove an entropy floor on recurrent
state (Proposition~\ref{prop:floor}), reparameterizing the known length-based $\Omega(n)$ bound by
the realized key$\rightarrow$value entropy $H(M)$; this is a modest, standard-technique result that
isolates the low-entropy regime as the only place a recurrent mechanism could win. \emph{Empirics:}
on a tested, reproducible harness of five architectures and five recall tasks, we reproduce the
recall wall and add the observation that \emph{content-based addressing}---not merely a growing
state---is what separates the transformer from the rest (a non-selective SSM and a content-blind
long convolution both fail). \emph{Mechanism:} our surprise-gated sparse-memory proposal (H2) is,
so far, a \textbf{negative result}: under recall-only supervision the gate does not sparsify, so it
neither saves state nor beats baselines.

We deliberately do not overclaim. The novel theorem is a corollary-grade reparameterization; the
mechanism is unconfirmed. What the project does establish is a rigorous, falsifiable, fully
reproducible pipeline---and an honest map of what works, what is already known, and what failed.

\paragraph{Future work, in priority order.}
(1) \textbf{Fix the gate signal:} add an auxiliary next-token loss so the backbone learns the
predictable structure and surprise becomes informative---the most direct test of whether H2's
state savings are achievable. (2) \textbf{Calibrate the structured task} so the upper-bound
transformer solves it, making the H2 comparison meaningful. (3) \textbf{Multi-seed and scale}
studies with error bars. (4) Tighten Proposition~\ref{prop:floor} in $\varepsilon$ and extend it
to segment-recurrent / chunked-attention hybrids that partially escape assumption~(A2). A
confirmed or cleanly-refuted H2 under (1)--(2) is the natural next milestone.
